\documentclass[a4paper,11pt,fleqn]{article}%fleqn pour aligner les equations à gauche
\usepackage{../../Styles/Cours}

\newcommand{\chnum}{Chapitre 13}
\newcommand{\chtitle}{La masse du Soleil}
\newcommand{\dtype}{TP2}

\begin{document}
\maketitle

\textbf{Objectif : Déterminer la masse du Soleil à partir d'observations}
\begin{figure}[h!]
	\begin{minipage}{.4\linewidth}
		\caption{Éphémérides de la Terre}
		\begin{tabular}{|>{\centering}m{.25\linewidth}|>{\centering}m{.25\linewidth}|>{\centering\arraybackslash}m{.25\linewidth}|}
			\hline Date $t$ (j) depuis le 01/01/2020 & Abscisse $x$ (u.a.) & Ordonnée $y$ (u.a.)\\
			\hline 0 & 0,0175 & 0,9998\\
			\hline 2 & -0,0170 & 0,9999\\
			\hline 4 & -0,0513 & 0,9987\\
			\hline 6 & -0,0857 & 0,9963\\
			\hline 8 & -0,1199 & 0,9928\\
			\hline 10 & -0,1540 & 0,9881\\
			\hline 12 & -0,1879 & 0,9822\\
			\hline
			\end{tabular}\\
			\textbf{Donnée :} 1 u.a. = \SI{1,5E8}{\kilo\meter}
	\end{minipage}\hspace{.02\linewidth}
	\begin{minipage}{.58\linewidth}
			 \caption{Code Python}
			\begin{lstlisting}[language = Python]
			import numpy as np
			import matplotlib.pyplot as plt
			t = []
			x = []
			y = []
			plt.xlabel('Coordonnees x')
			plt.ylabel('Coordonnees y')
			plt.title('Trajectoire de la Terre')
			plt.scatter(x, y, marker = '+')
			i = 1
			while i < len(t)-1 :
			  alpha = np.arctan2(y[i], x[i]) - np.arctan2(y[i-1], x[i-1])
			  r0 = np.sqrt(x[i]**2 + y[i]**2)
			  r1 = np.sqrt(x[i-1]**2 + y[i-1]**2)
			  A = r0*r1*np.sin(alpha)/2
			  plt.fill([x[i], x[i-1], 0], [y[i], y[i-1], 0], label='A = ' + "%.2e"%A + ' u.a.**2')
			  i += 2
			plt.legend(loc='center right')
			plt.show()
			\end{lstlisting}
	\end{minipage}
\end{figure}

\begin{figure}[h!]
	\caption{ Caractéristiques planétaires}
	\begin{tabular}{|>{\centering}m{2.5cm}|>{\centering}m{2cm}|>{\centering}m{2cm}|>{\centering}m{2cm}|>{\centering}m{2cm}|>{\centering}m{2cm}|>{\centering\arraybackslash}m{2cm}|}
		\hline Planète & Mercure & Vénus & Terre & Mars & Jupiter & Saturne\\
		\hline Demi-grand axe $a$ (u.a.) & \num{0,39} & \num{0,72} & \num{1,00} & \num{1,52} & \num{5,20} & \num{9,52}\\
		\hline Période de révolution $T$ (j) & \num{87,9} & \num{224,7} & \num{365,25} & \num{687} & \num{4331} & \num{1075}\\
		\hline
		\end{tabular}
\end{figure}

\textbf{Questions :}
\begin{enumerate}
	\item Préciser que est le référentiel choisi pour étudier le mouvement des planètes. En utilisant les éphémérides, on peut vérifier la 2$^\text{e}$ loi de Kepler à l'aide d'un programme Python qui calcule la surface balayée par une planète pendant des durées identiques.
	\item Dans l'extrait du code Python proposé, expliquer comment se réalise le calcul d'une surface entre deux positions.
	\item Tracer dans un tableau ou avec le langage Python la courbe $T^2=f(a^3)$.
	\item Justifier que le tracé $T^2=f(a^3)$ est conforme à la relation suivante :
	$$\dfrac{T^2}{a^3}=\dfrac{4 \pi ^2}{G \cdot M_S}$$
	\item En déduire l'expression de la masse du Soleil $M_S$
	\item Proposer alors un protocole pour déterminer la masse du Soleil $M_S$ à partir de la courbe $T^2=f(a^3)$. Le faire valider par le professeur avant de le réaliser.
\end{enumerate}
\end{document}